
Este capítulo desarrolló la arquitectura híbrida que integra el modelo mecanicista DEB del \cref{chap:cap3} con un estimador externo de biomasa basado en Procesos Gaussianos, resolviendo el problema de estados parcialmente observables que limitaba la aplicación operativa del modelo puro. A continuación se sintetizan los aportes principales y se establece el puente hacia el siguiente capítulo de la tesis.

\paragraph{Estimación de biomasa mediante GP.}
Se formuló y validó un estimador de biomasa húmeda total $\hat{B}(t)$ que opera como regresor probabilístico no paramétrico, consumiendo variables operativas medidas en tiempo real ($C(t)$, $T$, HR, $t$) y entrenado con cosechas destructivas espaciadas cada \num{5} días. La estrategia de validación \emph{leave-one-replica-out} demostró que el GP alcanza un MAPE inferior al \num{15}\% en la predicción de biomasa, con un coeficiente de determinación $R^{2} \geq \num{0.70}$, cumpliendo los requisitos de parsimonia y generalización exigidos por la escasez de datos experimentales ($N < \num{100}$ observaciones por campaña). La propagación de la varianza predictiva $\sigma^{2}_{\hat{B}}(t)$ al modelo DEB, mediante linealización del Jacobiano o simulación Monte Carlo, permite que la predicción de emisiones incluya un intervalo de confianza derivado de la incertidumbre en la variable de estado no observable.

\paragraph{Calibración en dos fases.}
La estrategia de calibración separó explícitamente los parámetros \emph{semilla} —$\kappa$, $m$, $a_{\max}$, $Y_B$, $Y_L$ fijos desde \textcite{eriksenDynamicModellingFeed2022}— de los parámetros \emph{libres} —$\tau_h$, $\beta_0$, $\rho_{\mathrm{conv}}$, $T_A$, $T_{\mathrm{opt}}$, $T_H$— ajustados por optimización restringida (L-BFGS-B con gradientes por diferencias finitas centradas). Esta división preservó la interpretabilidad biofísica del núcleo mecanicista mientras adaptaba el modelo a las condiciones locales del reactor. La función de pérdida, definida sobre la discrepancia de pendientes de emisión en ventanas cerradas de \num{30} min, evitó el sesgo de saturación que afectaría a una comparación de concentraciones acumuladas.

\paragraph{Confrontación modelo--sensor y diagnóstico con \ce{CH4}.}
El sistema híbrido predice la tasa de emisión de \ce{CO2} con un MAPE de \num{8.3}\% y un $R^{2} = \num{0.78}$, superando ampliamente el modelo lineal estático de referencia ($\rho_{\mathrm{RMSE}} = \num{0.39}$). El metano se mantuvo fuera de las ecuaciones de estado del DEB —conforme al supuesto de aerobiosis estricta— pero operó como canal paralelo de diagnóstico mediante un clasificador binario con umbral de \num{100} ppm. Durante la validación experimental, el detector alcanzó un recall del \num{100}\% en la identificación de eventos de anaerobiosis confirmados, validando la hipótesis H4 y demostrando utilidad como barrera de seguridad ambiental.

\paragraph{Inventario dinámico de carbono.}
La integración modelo--datos generó un inventario dinámico de carbono emitido por ciclo, cuantificando una discrepancia media acumulada $\bar{\Delta}_{\mathrm{Total}} = \num{+5.1}$ g C (\num{3.6}\% del total medido). Este residuo sistemático, atribuible a la respiración microbiana del sustrato no capturada por el DEB larval, es de magnitud reducida y permite su corrección mediante un factor de ajuste en aplicaciones operativas futuras.

\paragraph{Limitaciones y delimitación.}
Se reconoce que la arquitectura híbrida DEB+GP fue validada a escala de laboratorio (mesocosmos de \num{14} días, \num{6} réplicas). La transferencia a escala piloto o industrial exigirá: (i) recalibración del kernel del GP ante nuevos rangos de densidad larval y tipos de sustrato; (ii) adaptación de las restricciones de caja de los parámetros libres si la geometría del reactor cambia sustancialmente; y (iii) validación del desempeño del sensor NDIR en condiciones de humedad extrema prolongada (HR $> \num{90}$\%). No obstante, la estructura modular del sistema —separación explícita entre capa de inferencia (GP), capa mecanicista (DEB) y capa de diagnóstico (\ce{CH4})— facilita dichas adaptaciones sin requerir una reformulación del marco teórico.

\paragraph{Puente al Capítulo 5.}
Con la arquitectura híbrida validada y los parámetros calibrados, el siguiente capítulo traslada el sistema desde el dominio de la predicción hacia el dominio de la acción. Se desarrollará la plataforma de soporte a la decisión operativa, que integra: (i) el semáforo de riesgo ambiental (verde/ámbar/rojo) alimentado por los residuos $\delta(t_k)$ y el indicador $\mathcal{A}(t)$; (ii) el protocolo de respuesta ante eventos de anaerobiosis; y (iii) la generación automatizada de reportes MRV con trazabilidad de timestamp y sellado de datos, orientados a la verificación externa del desempeño ambiental del proceso de bioconversión. El Capítulo 5, en consecuencia, completa la transición del sistema híbrido desde una herramienta de estimación hacia una herramienta de gestión ambiental operativa.
